\centerline{\Huge \bf  Тренировочная Олимпиада 2}
\centerline{\Huge \bf 7 класс}
\centerline{\Huge \bf Уровень: Регион Эйлера. День 1.}

\begin{flushright}
    \scriptsize{Глава 2. Зверь.}
\end{flushright}


\problem{1} 
Клетки доски $16\times 16$ покрасили в один из двух цветов: красный или синий. Оказалось, что в любом трёхклеточном уголке красных клеток больше, чем синих. Найдите наибольшее возможное число синих клеток.
\begin{flushright}
    \textit{\small{(Гасанов Р.В.)}}
\end{flushright}

\problem{2} 
 Докажите, что при $x, y \geq \dfrac{1}{2}$ выполняется неравенство:
 
\[ x^2y^2 + 2(x+y) \geq 4xy + 1 \]
\begin{flushright}
    \textit{\small{(Гасанов Р.В.)}}
\end{flushright}

\problem{3} 
Мишень имеет форму правильного $6-$угольника со стороной 1 метр. Техасский бандит Дикарь Герман произвёл $7$ выстрелов по мишени со своего револьвера. Докажите, что  найдутся две пробоины расстояние между которыми меньше метра.
\begin{flushright}
    \textit{\small{(Гасанов Р.В.)}}
\end{flushright}

\problem{4}
Мияска рисовала на холсте свой очередной шедевр и случайно разлила краски на лежащую рядом шахматную доску. Некоторые \textit{покрашенные} клетки сразу высохли, а некоторые клетки остались \textit{липкими}. Один из двух ферзей тоже сразу покрасился, а другой остался липким. \textit{Покрашенный} ферзь может ходить по покрашенным клеткам, а \textit{липкий} ферзь только по липким (оба могут ходить по своим клеткам по несколько раз). Докажите, что какой-то из ферзей может пройтись по всем своим клеткам.
\begin{flushright}
    \textit{\small{(Гасанов Р.В.)}}
\end{flushright}

\problem{5}
Найдите все натуральные числа, которые нельзя представить в виде $\dfrac{m}{n} + \dfrac{m+1}{n+1}$, где $m, n$ натуральные числа.

\begin{flushright}
    \textit{\small{(Гасанов Р.В.)}}
\end{flushright}